\documentclass[aspectratio=169,10pt]{beamer}

% Shared Beamer setup for Fall 2026 course slides.
\mode<presentation>{
  \usetheme{Madrid}
  \usecolortheme{default}
}

\definecolor{CalPolyGreen}{HTML}{154734}
\definecolor{CalPolyGold}{HTML}{BD8B13}
\definecolor{DeckInk}{HTML}{1F2933}
\definecolor{DeckMuted}{HTML}{5F6B7A}

\setbeamercolor{structure}{fg=CalPolyGreen}
\setbeamercolor{palette primary}{bg=CalPolyGreen,fg=white}
\setbeamercolor{palette secondary}{bg=DeckInk,fg=white}
\setbeamercolor{palette tertiary}{bg=CalPolyGold,fg=DeckInk}
\setbeamercolor{frametitle}{bg=CalPolyGreen,fg=white}
\setbeamercolor{title}{fg=CalPolyGreen}
\setbeamercolor{normal text}{fg=DeckInk,bg=white}
\setbeamercolor{block title}{bg=CalPolyGreen,fg=white}
\setbeamercolor{block body}{bg=black!3,fg=DeckInk}

\setbeamertemplate{navigation symbols}{}
\setbeamertemplate{footline}[frame number]
\setbeamertemplate{itemize item}{\color{CalPolyGold}$\blacktriangleright$}
\setbeamertemplate{itemize subitem}{\color{CalPolyGreen}$\bullet$}

\newcommand{\CourseNumber}{Course}
\newcommand{\CourseTitle}{Course Title}
\newcommand{\LectureNumber}{0}
\newcommand{\LectureTitle}{Lecture Title}
\newcommand{\LectureDate}{Date}
\newcommand{\CourseUnit}{Unit}

\newcommand{\coursenumber}[1]{\renewcommand{\CourseNumber}{#1}}
\newcommand{\coursetitle}[1]{\renewcommand{\CourseTitle}{#1}}
\newcommand{\lecturenumber}[1]{\renewcommand{\LectureNumber}{#1}}
\newcommand{\lecturetitle}[1]{\renewcommand{\LectureTitle}{#1}}
\newcommand{\lecturedate}[1]{\renewcommand{\LectureDate}{#1}}
\newcommand{\courseunit}[1]{\renewcommand{\CourseUnit}{#1}}

\author{Kirk Duran}
\institute{California Polytechnic State University, San Luis Obispo}
\date{\LectureDate}

\newcommand{\makedecktitle}{
  \begin{frame}[plain]
    \vspace{0.25in}
    {\usebeamerfont{title}\color{CalPolyGreen}\LectureTitle\par}
    \vspace{0.18in}
    {\large \CourseNumber{}: \CourseTitle\par}
    \vspace{0.08in}
    {\large Lecture \LectureNumber{} -- \LectureDate\par}
    \vfill
    {\small Kirk Duran \textbar{} Cal Poly, San Luis Obispo\par}
  \end{frame}
}

\newcommand{\placeholdernote}{
  \begin{block}{Placeholder}
    Replace these bullets with the final lecture material, examples, and in-class activities.
  \end{block}
}

\AtBeginSection[]{
  \begin{frame}{Roadmap}
    \tableofcontents[currentsection]
  \end{frame}
}


\usepackage{tikz}

\graphicspath{{../assets/lecture-10/}}

% If an image file is not yet available, skip the visual slot.
\newcommand{\lectureimage}[2][1.75in]{%
  \IfFileExists{../assets/lecture-10/#2}{%
    \includegraphics[width=\linewidth,height=#1,keepaspectratio]{#2}%
  }{}%
}

% Explicit placeholder for visuals/examples that will be added later.
\newcommand{\lectureplaceholder}[2][2.35in]{%
  \begin{center}
    \fbox{%
      \begin{minipage}[c][#1][c]{0.90\linewidth}
        \centering
        \small\textit{#2}
      \end{minipage}%
    }
  \end{center}%
}

\setbeamersize{text margin left=0.42in,text margin right=0.42in}

\coursenumber{CS 4880}
\coursetitle{Artificial Intelligence}
\lecturenumber{10}
\lecturetitle{Monte Carlo Tree Search}
\lecturedate{September 22, 2026}
\courseunit{Search and Problem Solving}

\begin{document}

% -----------------------------------------------------------------------------
% Source slide 1 -> New slide 1
% -----------------------------------------------------------------------------
\makedecktitle

% -----------------------------------------------------------------------------
% Source slide 2 -> New slide 2
% -----------------------------------------------------------------------------
\begin{frame}[shrink=18]{Beyond Deep Blue: When Full-Width Search Breaks}
  \begin{columns}[T,onlytextwidth]
    \begin{column}{0.55\textwidth}
      Minimax and alpha--beta pruning work well when the game tree is manageable and a strong evaluation function is available.
      \begin{itemize}
        \item Chess has a large tree, but tactical structure and move ordering make pruning effective.
        \item Some games have much larger branching factors or unreliable handcrafted evaluation functions.
        \item Randomness or hidden information can further multiply the number of outcomes that must be considered.
      \end{itemize}
      \begin{block}{A different strategy}
        Instead of expanding every relevant continuation, can we \textbf{sample possible futures} and concentrate computation where it matters?
      \end{block}
    \end{column}
    \begin{column}{0.39\textwidth}
      \lectureplaceholder[2.55in]{Visual placeholder: chess tree transitioning into a much wider Go-like tree, with only a small sampled region highlighted.}
    \end{column}
  \end{columns}
  \begin{keyidea}
    Monte Carlo search trades exhaustive enumeration for statistical estimates built from repeated simulations.
  \end{keyidea}
\end{frame}

\section{Why Sampling?}

% -----------------------------------------------------------------------------
% Source slide 3 -> New slide 3
% -----------------------------------------------------------------------------
\begin{frame}[shrink=18]{Why Go Challenged Classical Game Search}
  \begin{columns}[T,onlytextwidth]
    \begin{column}{0.54\textwidth}
      Go is difficult for full-width adversarial search because its tree is both \textbf{wide} and \textbf{deep}.
      \begin{itemize}
        \item A typical early position may offer hundreds of legal moves.
        \item The number of legal board configurations is on the order of $10^{170}$.
        \item Local moves can have global strategic consequences many turns later.
        \item Exhaustive search to terminal states is completely infeasible.
      \end{itemize}
      \begin{block}{Consequence}
        Even dramatic pruning leaves far too many continuations to explore uniformly.
      \end{block}
    \end{column}
    \begin{column}{0.40\textwidth}
      \lectureplaceholder[2.65in]{Visual placeholder: chess board beside Go board, with a narrow chess tree and dramatically wider Go tree.}
    \end{column}
  \end{columns}
\end{frame}

% -----------------------------------------------------------------------------
% Source slide 4 -> New slide 4
% -----------------------------------------------------------------------------
\begin{frame}[shrink=18]{Why Alpha--Beta Is Not Enough}
  \begin{columns}[T,onlytextwidth]
    \begin{column}{0.54\textwidth}
      Alpha--beta pruning is exact, but its efficiency depends heavily on the structure of the search.
      \begin{itemize}
        \item Good move ordering is required for aggressive pruning.
        \item A depth limit still requires an informative evaluation function.
        \item In Go, many positions are strategically ambiguous until much later.
        \item Searching every move to the same depth wastes effort on weak or irrelevant branches.
      \end{itemize}
      \begin{block}{Shift in viewpoint}
        Rather than ask ``How deep can we search every move?'', ask ``Which moves deserve more samples?''
      \end{block}
    \end{column}
    \begin{column}{0.40\textwidth}
      \lectureplaceholder[2.60in]{Visual placeholder: a full-width alpha--beta tree contrasted with an uneven tree that allocates more depth to promising branches.}
    \end{column}
  \end{columns}
  \begin{keyidea}
    MCTS is a selective search method: computation is allocated adaptively rather than uniformly by depth.
  \end{keyidea}
\end{frame}

% -----------------------------------------------------------------------------
% Source slide 5 -> New slide 5
% -----------------------------------------------------------------------------
\begin{frame}[shrink=18]{Reminder: Deterministic vs. Stochastic Games}
  \begin{columns}[T,onlytextwidth]
    \begin{column}{0.47\textwidth}
      \begin{block}{Deterministic adversarial games}
        \begin{itemize}
          \item Actions have deterministic outcomes.
          \item Minimax models an opponent choosing the worst response for MAX.
          \item Example: chess under perfect-information assumptions.
        \end{itemize}
      \end{block}
    \end{column}
    \begin{column}{0.47\textwidth}
      \begin{block}{Stochastic games}
        \begin{itemize}
          \item Some transitions are random.
          \item Chance nodes represent probability distributions over outcomes.
          \item Examples: dice, shuffled cards, random events.
        \end{itemize}
      \end{block}
    \end{column}
  \end{columns}
  \vspace{0.08in}
  \lectureplaceholder[1.35in]{Visual placeholder: chess piece on the left; dice/cards on the right; arrows showing deterministic choice versus probabilistic outcome.}
\end{frame}

% -----------------------------------------------------------------------------
% Source slide 6 -> New slide 6
% -----------------------------------------------------------------------------
\begin{frame}[shrink=18]{Expectiminimax Recap}
  \begin{cpdefinition}
    \textbf{Expectiminimax} extends minimax with \textbf{chance nodes} whose values are probability-weighted expectations.
  \end{cpdefinition}
  \begin{columns}[T,onlytextwidth]
    \begin{column}{0.53\textwidth}
      \begin{itemize}
        \item MAX node: choose the child with maximum value.
        \item MIN node: choose the child with minimum value.
        \item CHANCE node: compute
      \end{itemize}
      \[
        V(s)=\sum_i P(s_i\mid s)\,V(s_i).
      \]
      \begin{block}{Assumption}
        The probability model for each chance outcome is known and can be enumerated.
      \end{block}
    \end{column}
    \begin{column}{0.41\textwidth}
      \lectureplaceholder[2.35in]{Example placeholder: MAX node above CHANCE nodes, then MIN nodes and terminal utilities. Label chance edges with probabilities.}
    \end{column}
  \end{columns}
\end{frame}

% -----------------------------------------------------------------------------
% Source slide 7 -> New slide 7
% -----------------------------------------------------------------------------
\begin{frame}[shrink=18]{Expectiminimax: One Worked Backup}
  \begin{columns}[T,onlytextwidth]
    \begin{column}{0.54\textwidth}
      Suppose a chance event selects between two MIN states with equal probability.
      \[
        \min(3,5)=3, \qquad \min(2,4)=2.
      \]
      Therefore the chance node has value
      \[
        0.5(3)+0.5(2)=2.5.
      \]
      \begin{itemize}
        \item Exact expectation requires both outcomes.
        \item More outcomes mean more branches to enumerate.
      \end{itemize}
    \end{column}
    \begin{column}{0.40\textwidth}
      \lectureplaceholder[2.45in]{Example placeholder: chance node with two $0.5$ branches to MIN nodes valued $3$ and $2$, backing up expected value $2.5$.}
    \end{column}
  \end{columns}
  \begin{keyidea}
    Expectation is easy to compute when the outcome set is small; it becomes expensive when uncertainty creates many possible futures.
  \end{keyidea}
\end{frame}

% -----------------------------------------------------------------------------
% Source slide 8 -> New slide 8
% -----------------------------------------------------------------------------
\begin{frame}[shrink=18]{Stochastic Games Multiply the Branching Factor}
  \begin{columns}[T,onlytextwidth]
    \begin{column}{0.52\textwidth}
      Random events can make a game tree much wider than the action set alone suggests.
      \begin{itemize}
        \item Backgammon combines legal moves with dice outcomes.
        \item Card games combine actions with hidden deals and later draws.
        \item Each additional chance outcome creates another continuation to evaluate.
      \end{itemize}
      \begin{block}{Effective branching}
        A player branch followed by a chance branch behaves roughly like the product of their branching factors.
      \end{block}
    \end{column}
    \begin{column}{0.42\textwidth}
      \lectureplaceholder[2.55in]{Visual placeholder: backgammon board with a decision node followed by a large fan-out of dice outcomes.}
    \end{column}
  \end{columns}
\end{frame}

% -----------------------------------------------------------------------------
% Source slide 9 -> New slide 9
% -----------------------------------------------------------------------------
\begin{frame}[shrink=18]{The Limitation of Exhaustive Expectation}
  \begin{columns}[T,onlytextwidth]
    \begin{column}{0.52\textwidth}
      Exact expectiminimax inherits the exponential growth of minimax and adds chance branching.
      \begin{itemize}
        \item Every modeled outcome must be expanded.
        \item Low-probability outcomes still consume search effort.
        \item Exact probabilities may be expensive or unknown.
        \item A fixed-depth tree does not naturally provide a useful answer if interrupted early.
      \end{itemize}
      \begin{block}{Question}
        Can we estimate an expectation from a \textbf{representative sample} instead of enumerating every outcome?
      \end{block}
    \end{column}
    \begin{column}{0.42\textwidth}
      \lectureplaceholder[2.50in]{Visual placeholder: exhaustive expectimax tree with many faint chance branches; only several sampled trajectories emphasized.}
    \end{column}
  \end{columns}
\end{frame}

% -----------------------------------------------------------------------------
% Source slide 10 -> New slide 10
% -----------------------------------------------------------------------------
\begin{frame}[shrink=18]{From Enumeration to Sampling}
  \begin{columns}[T,onlytextwidth]
    \begin{column}{0.52\textwidth}
      Consider again a chance node with two equally likely outcomes worth $3$ and $2$.
      \begin{itemize}
        \item Exact value: $0.5(3)+0.5(2)=2.5$.
        \item Sample one outcome: estimate is either $3$ or $2$.
        \item Sample repeatedly and average: the estimate approaches $2.5$.
      \end{itemize}
      \begin{block}{Monte Carlo principle}
        Approximate an expectation using empirical averages from random samples.
      \end{block}
    \end{column}
    \begin{column}{0.42\textwidth}
      \lectureplaceholder[2.50in]{Example placeholder: same expectiminimax tree, but show repeated sampled paths and a running average converging toward $2.5$.}
    \end{column}
  \end{columns}
  \begin{keyidea}
    Sampling replaces complete knowledge of the tree with progressively better statistical estimates.
  \end{keyidea}
\end{frame}

\section{Bandits and Tree Selection}

% -----------------------------------------------------------------------------
% Source slide 11 -> New slide 11
% -----------------------------------------------------------------------------
\begin{frame}[shrink=18]{Monte Carlo Tree Search: High-Level Idea}
  \begin{cpdefinition}
    \textbf{Monte Carlo Tree Search (MCTS)} incrementally builds a partial search tree by repeatedly selecting, expanding, simulating, and backing up sampled outcomes.
  \end{cpdefinition}
  \begin{columns}[T,onlytextwidth]
    \begin{column}{0.54\textwidth}
      \begin{itemize}
        \item Promising actions receive more simulations.
        \item Poorly performing actions gradually receive less attention.
        \item The tree becomes deeper where evidence suggests depth is valuable.
        \item Search can be stopped at almost any time and still return a current best move.
      \end{itemize}
      \begin{block}{Modern variants}
        MCTS can use random rollouts, heuristic rollouts, or learned policy/value models.
      \end{block}
    \end{column}
    \begin{column}{0.40\textwidth}
      \lectureplaceholder[2.15in]{Visual placeholder: asymmetric MCTS tree with a few heavily visited branches and many lightly explored alternatives.}
    \end{column}
  \end{columns}
\end{frame}

% -----------------------------------------------------------------------------
% Source slide 12 -> New slide 12
% -----------------------------------------------------------------------------
\begin{frame}[shrink=18]{The Multi-Armed Bandit Problem}
  \begin{columns}[T,onlytextwidth]
    \begin{column}{0.54\textwidth}
      Imagine $k$ slot-machine arms with unknown reward distributions.
      \begin{itemize}
        \item Pulling an arm reveals a reward sample.
        \item We want high total reward over repeated pulls.
        \item The best-known arm may not truly be the best arm.
      \end{itemize}
      \begin{block}{The central tension}
        \textbf{Exploit} actions that currently look strong, or \textbf{explore} actions whose value is still uncertain.
      \end{block}
    \end{column}
    \begin{column}{0.40\textwidth}
      \lectureplaceholder[2.45in]{Visual placeholder: three slot-machine arms labeled A, B, C with unknown reward distributions and a question mark over which to choose.}
    \end{column}
  \end{columns}
  \begin{keyidea}
    Every interior MCTS node creates a local bandit problem over its available actions.
  \end{keyidea}
\end{frame}

% -----------------------------------------------------------------------------
% Source slide 13 -> New slide 13
% -----------------------------------------------------------------------------
\begin{frame}[shrink=18]{Bandit Example: Which Arm Next?}
  \begin{columns}[T,onlytextwidth]
    \begin{column}{0.50\textwidth}
      Suppose we have observed:
      \begin{center}
      \renewcommand{\arraystretch}{1.25}
      \begin{tabular}{c|c|c|c}
        \textbf{Arm} & \textbf{Visits} & \textbf{Total reward} & \textbf{Mean} \\ \hline
        A & 3 & 2 & $0.67$ \\
        B & 2 & 1 & $0.50$ \\
        C & 1 & 1 & $1.00$
      \end{tabular}
      \end{center}
      \begin{itemize}
        \item C has the highest observed mean.
        \item C is also the least tested.
        \item A has more evidence, but a lower mean.
      \end{itemize}
    \end{column}
    \begin{column}{0.44\textwidth}
      \lectureplaceholder[2.45in]{Example placeholder: three arms annotated with visit counts and rewards; ask ``Which arm should be pulled next?''}
    \end{column}
  \end{columns}
\end{frame}

% -----------------------------------------------------------------------------
% Source slide 14 -> New slide 14
% -----------------------------------------------------------------------------
\begin{frame}[shrink=16]{Upper Confidence Bounds}
  \begin{cpdefinition}
    UCB1 scores each action using its estimated reward plus an uncertainty bonus.
  \end{cpdefinition}
  \[
    \operatorname{UCB1}(i)
    = \underbrace{\frac{Q_i}{N_i}}_{\text{exploitation}}
      + C\underbrace{\sqrt{\frac{\ln N}{N_i}}}_{\text{exploration}}.
  \]
  \begin{itemize}
    \item $Q_i$: cumulative reward observed for action $i$.
    \item $N_i$: number of times action $i$ has been selected.
    \item $N$: total selections at the parent.
    \item $C>0$: exploration constant controlling the tradeoff.
  \end{itemize}
  \begin{keyidea}
    Frequently successful actions score well because of their mean reward; rarely tested actions score well because uncertainty remains high.
  \end{keyidea}
\end{frame}

% -----------------------------------------------------------------------------
% Source slide 15 -> New slide 15
% -----------------------------------------------------------------------------
\begin{frame}[shrink=18]{UCB1 Intuition: Revisit or Explore?}
  \begin{columns}[T,onlytextwidth]
    \begin{column}{0.47\textwidth}
      \begin{block}{Exploitation}
        Return to an option that has already produced strong average rewards.
        \begin{itemize}
          \item evidence of quality;
          \item lower uncertainty;
          \item immediate payoff.
        \end{itemize}
      \end{block}
    \end{column}
    \begin{column}{0.47\textwidth}
      \begin{block}{Exploration}
        Try an option that has received fewer samples.
        \begin{itemize}
          \item information gain;
          \item higher uncertainty;
          \item possibility of discovering a better action.
        \end{itemize}
      \end{block}
    \end{column}
  \end{columns}
  \vspace{0.08in}
  \lectureplaceholder[1.20in]{Analogy placeholder: familiar food stall with a proven rating versus a new stall with little evidence but uncertain potential.}
\end{frame}

% -----------------------------------------------------------------------------
% Source slide 16 -> New slide 16
% -----------------------------------------------------------------------------
\begin{frame}[shrink=14]{UCB1 Worked Example}
  Using the bandit statistics from the previous slide, let $N=6$ and $C=\sqrt{2}$.
  \begin{align*}
    \operatorname{UCB1}(A)
      &= \frac{2}{3}+\sqrt{2}\sqrt{\frac{\ln 6}{3}}
      \approx 1.76,\\
    \operatorname{UCB1}(B)
      &= \frac{1}{2}+\sqrt{2}\sqrt{\frac{\ln 6}{2}}
      \approx 1.84,\\
    \operatorname{UCB1}(C)
      &= 1+\sqrt{2}\sqrt{\ln 6}
      \approx 2.89.
  \end{align*}
  \begin{block}{Decision}
    Choose C. Its observed reward is strong, and its single visit gives it the largest uncertainty bonus.
  \end{block}
  \begin{keyidea}
    UCB does not explore randomly; it explores where uncertainty is still plausibly valuable.
  \end{keyidea}
\end{frame}

% -----------------------------------------------------------------------------
% Source slide 17 -> New slide 17
% -----------------------------------------------------------------------------
\begin{frame}[shrink=18]{From UCB1 to UCT}
  \begin{columns}[T,onlytextwidth]
    \begin{column}{0.54\textwidth}
      MCTS applies the bandit idea recursively inside a search tree.
      \begin{itemize}
        \item Each state is a node.
        \item Each legal action is an arm.
        \item Simulated outcomes provide reward samples.
        \item The UCB-style rule used in trees is commonly called \textbf{UCT}: Upper Confidence bounds applied to Trees.
      \end{itemize}
      \begin{block}{Adversarial detail}
        Values must use a consistent player perspective. At opponent nodes, rewards are typically negated or stored from the player-to-move perspective.
      \end{block}
    \end{column}
    \begin{column}{0.40\textwidth}
      \lectureplaceholder[2.55in]{Visual placeholder: one search-tree node transformed into a local multi-armed-bandit decision over its child actions.}
    \end{column}
  \end{columns}
\end{frame}

\section{The MCTS Algorithm}

% -----------------------------------------------------------------------------
% Source slide 18 -> New slide 18
% -----------------------------------------------------------------------------
\begin{frame}[shrink=16]{The Four Phases of MCTS}
  \begin{columns}[T,onlytextwidth]
    \begin{column}{0.55\textwidth}
      \begin{enumerate}
        \item \textbf{Selection}: follow the tree from the root using UCT until reaching a node that is not fully expanded.
        \item \textbf{Expansion}: add one or more previously unvisited actions as new child nodes.
        \item \textbf{Simulation / Evaluation}: estimate the new node's outcome using a rollout or value model.
        \item \textbf{Backpropagation}: update reward statistics and visit counts along the selected path.
      \end{enumerate}
      \begin{block}{Repeat}
        Each iteration contributes one additional sample of evidence to the tree.
      \end{block}
    \end{column}
    \begin{column}{0.39\textwidth}
      \lectureplaceholder[2.75in]{Diagram placeholder: tree path labeled 1 Selection, new child labeled 2 Expansion, rollout trajectory labeled 3 Simulation, upward arrows labeled 4 Backpropagation.}
    \end{column}
  \end{columns}
\end{frame}

% -----------------------------------------------------------------------------
% Source slide 19 -> New slide 19
% -----------------------------------------------------------------------------
\begin{frame}[shrink=14]{MCTS Pseudocode}
  \begin{block}{Algorithmic structure}
  \ttfamily\scriptsize
  while budget remains:\\
  \quad node $\leftarrow$ root\\[0.03in]
  \quad // SELECTION\\
  \quad while node is fully expanded and nonterminal:\\
  \qquad node $\leftarrow$ child with best UCT score\\[0.03in]
  \quad // EXPANSION\\
  \quad if node is nonterminal:\\
  \qquad node $\leftarrow$ EXPAND(node)\\[0.03in]
  \quad // SIMULATION / EVALUATION\\
  \quad reward $\leftarrow$ ESTIMATE(node)\\[0.03in]
  \quad // BACKPROPAGATION\\
  \quad while node $\neq$ null:\\
  \qquad node.visits $\leftarrow$ node.visits $+1$\\
  \qquad node.value $\leftarrow$ node.value $+$ reward\\
  \qquad reward $\leftarrow$ TRANSFORM\_FOR\_PARENT(reward)\\
  \qquad node $\leftarrow$ node.parent
  \end{block}
  \begin{keyidea}
    The search tree is not generated in advance; it emerges from repeated evidence-gathering iterations.
  \end{keyidea}
\end{frame}

% -----------------------------------------------------------------------------
% Source slide 20 -> New slide 20
% -----------------------------------------------------------------------------
\begin{frame}[shrink=18]{MCTS from Scratch: Iteration 1}
  \begin{columns}[T,onlytextwidth]
    \begin{column}{0.52\textwidth}
      Begin with an unexpanded root $s_0$.
      \begin{enumerate}
        \item Expand left child $s_{10}$.
        \item Simulate from $s_{10}$.
        \item Suppose the rollout is a win for the root player: reward $1.0$.
        \item Backpropagate the result.
      \end{enumerate}
      \[
        N(s_{10})=1,\; Q(s_{10})=1.0
      \]
      \[
        N(s_0)=1,\; Q(s_0)=1.0.
      \]
      \begin{block}{Trace convention}
        To isolate the statistics, this toy trace stores rewards from the root player's perspective.
      \end{block}
    \end{column}
    \begin{column}{0.42\textwidth}
      \lectureplaceholder[2.65in]{Example placeholder: root $s_0$ with one child $s_{10}$; annotate both nodes with $N=1,Q=1.0$.}
    \end{column}
  \end{columns}
\end{frame}

% -----------------------------------------------------------------------------
% Source slide 21 -> New slide 21
% -----------------------------------------------------------------------------
\begin{frame}[shrink=16]{MCTS from Scratch: Iteration 2}
  \begin{columns}[T,onlytextwidth]
    \begin{column}{0.55\textwidth}
      The right child has never been visited.
      \begin{itemize}
        \item Treat an unvisited action as having effectively infinite exploration priority.
        \item Expand $s_{11}$.
        \item Suppose its rollout is a draw: reward $0.5$.
      \end{itemize}
      After backpropagation:
      \[
        N(s_{10})=1,\;Q(s_{10})=1.0
      \]
      \[
        N(s_{11})=1,\;Q(s_{11})=0.5
      \]
      \[
        N(s_0)=2,\;Q(s_0)=1.5.
      \]
    \end{column}
    \begin{column}{0.39\textwidth}
      \lectureplaceholder[2.65in]{Example placeholder: root with two children, left labeled $N=1,Q=1.0$, right labeled $N=1,Q=0.5$, root $N=2,Q=1.5$.}
    \end{column}
  \end{columns}
\end{frame}

% -----------------------------------------------------------------------------
% Source slide 22 -> New slide 22
% -----------------------------------------------------------------------------
\begin{frame}[shrink=14]{MCTS from Scratch: Iteration 3}
  With $C=\sqrt{2}$ and $N(s_0)=2$:
  \begin{align*}
    \operatorname{UCT}(s_{10})
      &= 1.0 + \sqrt{2}\sqrt{\ln 2}
      \approx 2.18,\\
    \operatorname{UCT}(s_{11})
      &= 0.5 + \sqrt{2}\sqrt{\ln 2}
      \approx 1.68.
  \end{align*}
  Select $s_{10}$ and expand child $s_{20}$. Suppose the rollout is a loss: reward $0.0$.
  \begin{center}
  \renewcommand{\arraystretch}{1.20}
  \begin{tabular}{c|c|c}
    \textbf{Node} & $N$ & $Q$ \\ \hline
    $s_{20}$ & 1 & 0.0 \\
    $s_{10}$ & 2 & 1.0 \\
    $s_{11}$ & 1 & 0.5 \\
    $s_0$    & 3 & 1.5
  \end{tabular}
  \end{center}
  \begin{keyidea}
    A disappointing rollout lowers the empirical value of the branch without permanently eliminating it.
  \end{keyidea}
\end{frame}

% -----------------------------------------------------------------------------
% Source slide 23 -> New slide 23
% -----------------------------------------------------------------------------
\begin{frame}[shrink=13]{MCTS from Scratch: Iteration 4}
  At the root, both actions now have mean reward $0.5$, but the less-visited action receives a larger exploration bonus.
  \begin{align*}
    \operatorname{UCT}(s_{10})
      &= \frac{1.0}{2}+\sqrt{2}\sqrt{\frac{\ln 3}{2}}
      \approx 1.55,\\
    \operatorname{UCT}(s_{11})
      &= \frac{0.5}{1}+\sqrt{2}\sqrt{\ln 3}
      \approx 1.98.
  \end{align*}
  Select $s_{11}$, expand $s_{22}$, and suppose the rollout wins with reward $1.0$.
  \[
    N(s_{22})=1,Q(s_{22})=1.0;
    \quad N(s_{11})=2,Q(s_{11})=1.5;
    \quad N(s_0)=4,Q(s_0)=2.5.
  \]
  \begin{block}{What changed?}
    UCT redirected computation toward the less-certain branch, then updated its estimate using the new evidence.
  \end{block}
\end{frame}

% -----------------------------------------------------------------------------
% Source slide 24 -> New slide 24
% -----------------------------------------------------------------------------
\begin{frame}[shrink=18]{What a Mature MCTS Tree Looks Like}
  \begin{columns}[T,onlytextwidth]
    \begin{column}{0.53\textwidth}
      After many iterations, the search tree becomes highly asymmetric.
      \begin{itemize}
        \item Strong actions receive many visits.
        \item Weak actions may remain shallow.
        \item Deep search emerges only in promising regions.
        \item The best current move can be returned when the time budget expires.
      \end{itemize}
      \begin{block}{Typical final decision}
        Choose the root action with the largest visit count or strongest robust value estimate.
      \end{block}
    \end{column}
    \begin{column}{0.41\textwidth}
      \lectureplaceholder[2.65in]{Visual placeholder: mature asymmetric tree, with edge thickness proportional to visits and only a few deep branches.}
    \end{column}
  \end{columns}
  \begin{keyidea}
    MCTS spends computation according to evidence, not according to a uniform depth boundary.
  \end{keyidea}
\end{frame}

\section{Rollouts, Comparisons, and Extensions}

% -----------------------------------------------------------------------------
% Source slide 25 -> New slide 25
% -----------------------------------------------------------------------------
\begin{frame}[shrink=18]{The Risk of Random Rollouts}
  \begin{columns}[T,onlytextwidth]
    \begin{column}{0.54\textwidth}
      Uniformly random simulation is simple, but its estimates can be poor.
      \begin{itemize}
        \item \textbf{Variance}: a small number of lucky rollouts may distort the mean.
        \item \textbf{Bias}: random play may not resemble competent play.
        \item \textbf{Blind spots}: rare tactical sequences may almost never be sampled.
        \item \textbf{Waste}: obviously bad actions can consume rollout time.
      \end{itemize}
      \begin{block}{Statistical consequence}
        Noisy rollouts require more samples before value estimates become reliable.
      \end{block}
    \end{column}
    \begin{column}{0.40\textwidth}
      \lectureplaceholder[2.55in]{Visual placeholder: several rollout paths with widely varying outcomes; label variance, tactical blind spot, and wasted simulation.}
    \end{column}
  \end{columns}
\end{frame}

% -----------------------------------------------------------------------------
% Source slide 26 -> New slide 26
% -----------------------------------------------------------------------------
\begin{frame}[shrink=18]{Improving the Simulation Policy}
  \begin{columns}[T,onlytextwidth]
    \begin{column}{0.30\textwidth}
      \begin{block}{Heuristic rollouts}
        Bias simulations toward plausible actions using domain knowledge.
      \end{block}
    \end{column}
    \begin{column}{0.30\textwidth}
      \begin{block}{Policy models}
        Use a learned policy to sample stronger, more realistic continuations.
      \end{block}
    \end{column}
    \begin{column}{0.30\textwidth}
      \begin{block}{Leaf evaluation}
        Stop the rollout early and estimate the state with a heuristic or learned value model.
      \end{block}
    \end{column}
  \end{columns}
  \vspace{0.15in}
  \begin{itemize}
    \item Better simulations can reduce variance and improve sample efficiency.
    \item Stronger guidance introduces model bias, so search and evaluation must be balanced carefully.
  \end{itemize}
  \begin{keyidea}
    Modern MCTS often replaces long random rollouts with learned policies and value estimates.
  \end{keyidea}
\end{frame}

% -----------------------------------------------------------------------------
% Source slide 27 -> New slide 27
% -----------------------------------------------------------------------------
\begin{frame}[shrink=16]{MCTS vs. Alpha--Beta Search}
  \begin{center}
  \renewcommand{\arraystretch}{1.25}
  \begin{tabular}{p{0.22\textwidth}|p{0.34\textwidth}|p{0.34\textwidth}}
    & \textbf{Alpha--Beta} & \textbf{MCTS} \\ \hline
    Search shape & depth-oriented, full-width before pruning & asymmetric, sample-driven \\
    Value estimate & heuristic evaluation at cutoff & empirical returns and/or learned value \\
    Strength & deterministic games with strong evaluation and move ordering & enormous branching factors and anytime search \\
    Weakness & sensitive to branching factor and evaluation quality & noisy estimates; may need many simulations \\
    Decision budget & usually fixed depth / iterative deepening & naturally simulation- or time-bounded
  \end{tabular}
  \end{center}
  \begin{keyidea}
    Neither method dominates universally: the best search algorithm depends on game structure, available models, and computational budget.
  \end{keyidea}
\end{frame}

% -----------------------------------------------------------------------------
% Source slide 28 -> New slide 28
% -----------------------------------------------------------------------------
\begin{frame}[shrink=18]{Games with Partial Observability}
  \begin{columns}[T,onlytextwidth]
    \begin{column}{0.54\textwidth}
      In imperfect-information games, the true state is not fully observed.
      \begin{itemize}
        \item Poker: opponents' private cards are hidden.
        \item Stratego: opponent piece identities are initially hidden.
        \item Real-time strategy games may contain fog of war.
      \end{itemize}
      \begin{cpdefinition}
        A \textbf{belief state} represents uncertainty over possible underlying states consistent with the observations so far.
      \end{cpdefinition}
    \end{column}
    \begin{column}{0.40\textwidth}
      \lectureplaceholder[2.50in]{Visual placeholder: observed public state leading to several possible hidden worlds inside a belief-state cloud.}
    \end{column}
  \end{columns}
  \begin{keyidea}
    Search must now reason over information, not only over physical game states.
  \end{keyidea}
\end{frame}

% -----------------------------------------------------------------------------
% Source slide 29 -> New slide 29
% -----------------------------------------------------------------------------
\begin{frame}[shrink=18]{Belief Sampling and Determinization}
  \begin{columns}[T,onlytextwidth]
    \begin{column}{0.54\textwidth}
      A simple way to adapt MCTS is to sample possible hidden states.
      \begin{enumerate}
        \item Sample a complete world consistent with current observations.
        \item Run search as though that sampled world were the true state.
        \item Repeat across many sampled worlds.
        \item Aggregate the evidence for available actions.
      \end{enumerate}
      \begin{block}{Important caveat}
        Naive determinization can produce \emph{strategy fusion}: it may choose different actions in sampled worlds that the real agent cannot distinguish.
      \end{block}
    \end{column}
    \begin{column}{0.40\textwidth}
      \lectureplaceholder[2.55in]{Visual placeholder: one observed card-game state branching into several sampled hidden hands, each feeding a search tree.}
    \end{column}
  \end{columns}
\end{frame}

% -----------------------------------------------------------------------------
% Source slide 30 -> New slide 30
% -----------------------------------------------------------------------------
\begin{frame}[shrink=18]{MCTS in Poker-Like Games}
  \begin{columns}[T,onlytextwidth]
    \begin{column}{0.52\textwidth}
      Poker combines several complications at once:
      \begin{itemize}
        \item private information;
        \item stochastic card deals;
        \item adversarial betting decisions;
        \item opponent modeling and bluffing.
      \end{itemize}
      A practical MCTS-style system may:
      \begin{enumerate}
        \item sample hidden hands consistent with observations;
        \item simulate betting and future cards;
        \item back up utilities such as expected chip gain.
      \end{enumerate}
    \end{column}
    \begin{column}{0.42\textwidth}
      \lectureplaceholder[2.70in]{Example placeholder: poker table showing public cards, hidden opponent cards, sampled determinizations, and rollout branches for fold/call/raise.}
    \end{column}
  \end{columns}
  \begin{keyidea}
    Imperfect-information search requires care: information-set methods are preferable to pretending hidden state is known.
  \end{keyidea}
\end{frame}

% -----------------------------------------------------------------------------
% Source slide 31 -> New slide 31
% -----------------------------------------------------------------------------
\begin{frame}[shrink=18]{Worked / Video Example: Poker Decision}
  \begin{columns}[T,onlytextwidth]
    \begin{column}{0.52\textwidth}
      Suggested in-class walkthrough:
      \begin{itemize}
        \item identify the information available to the acting player;
        \item list candidate actions: fold, call, raise;
        \item sample several plausible opponent holdings;
        \item trace one simulation for each candidate action;
        \item compare estimated expected utility after repeated samples.
      \end{itemize}
      \begin{block}{Discussion question}
        What information would a naive determinization method accidentally give the agent that a real player does not possess?
      \end{block}
    \end{column}
    \begin{column}{0.42\textwidth}
      \lectureplaceholder[2.65in]{Media placeholder: short poker hand clip or sequence of screenshots for an MCTS / belief-sampling walkthrough.}
    \end{column}
  \end{columns}
\end{frame}

% -----------------------------------------------------------------------------
% Source slide 32 -> New slide 32
% -----------------------------------------------------------------------------
\begin{frame}[shrink=18]{Why Belief-State Search Is Expensive}
  \begin{columns}[T,onlytextwidth]
    \begin{column}{0.47\textwidth}
      \begin{block}{Many possible worlds}
        The number of hidden configurations may grow combinatorially with unknown cards, pieces, or states.
      \end{block}
      \begin{block}{Belief updates}
        New observations require probabilities over possible worlds to be revised.
      \end{block}
    \end{column}
    \begin{column}{0.47\textwidth}
      \begin{block}{Planning over information}
        The agent must choose actions that work across states it cannot distinguish.
      \end{block}
      \begin{block}{Opponent uncertainty}
        Good play may require reasoning about what the opponent knows and how the opponent acts.
      \end{block}
    \end{column}
  \end{columns}
  \begin{keyidea}
    Partial observability adds a second search problem: reasoning about the state of knowledge itself.
  \end{keyidea}
\end{frame}

% -----------------------------------------------------------------------------
% Source slide 33 -> New slide 33
% -----------------------------------------------------------------------------
\begin{frame}[shrink=16]{AlphaGo: MCTS Guided by Neural Networks}
  \begin{columns}[T,onlytextwidth]
    \begin{column}{0.54\textwidth}
      AlphaGo demonstrated that MCTS becomes far stronger when search is guided by learned models.
      \begin{itemize}
        \item A \textbf{policy network} prioritizes promising moves during tree search.
        \item A \textbf{value network} estimates the probability of winning from a position.
        \item Search statistics refine these learned estimates through lookahead.
        \item The result is a focused tree that combines prior knowledge with online simulation.
      \end{itemize}
      \begin{block}{Conceptual synthesis}
        Learning proposes and evaluates; MCTS allocates computation and tests consequences.
      \end{block}
    \end{column}
    \begin{column}{0.40\textwidth}
      \lectureplaceholder[2.65in]{Visual placeholder: Go board feeding policy network and value network into an MCTS tree; optionally include AlphaGo documentary image.}
    \end{column}
  \end{columns}
\end{frame}

% -----------------------------------------------------------------------------
% Source slide 34 -> New slide 34
% -----------------------------------------------------------------------------
\begin{frame}[shrink=18]{Takeaways}
  \begin{itemize}
    \item Full-width minimax and expectiminimax become impractical when branching factors are enormous.
    \item Monte Carlo methods estimate values from samples rather than enumerating every possible future.
    \item UCB1 and UCT balance \textbf{exploitation} of strong actions with \textbf{exploration} of uncertain actions.
    \item MCTS repeatedly performs \textbf{selection, expansion, simulation/evaluation, and backpropagation}.
    \item The resulting tree grows asymmetrically toward regions that appear promising.
    \item Rollout quality, value estimation, and information handling strongly affect performance.
    \item Neural guidance can make MCTS dramatically more sample-efficient, as demonstrated by AlphaGo.
    \item Next lecture: \textbf{Knowledge-Based Agents}.
  \end{itemize}
  \begin{keyidea}
    MCTS turns search into sequential evidence gathering: spend simulations where they are most likely to improve the decision.
  \end{keyidea}
\end{frame}

\end{document}
